\subsection{July}

\subsubsection{July 1st}
Today I learned that the Zeckendorf decomposition actually defines the Fibonacci sequence. To set out indexing properly, we fix our definition of the Fibonacci numbers.
\begin{definition}
    The Fibonacci numbers $\{F_k\}_{k=0}^\infty$ are defined with $F_0=1,$ $F_1=2,$ and
\end{definition}
This choice of indexing assures that all terms are indices are positive integers, and the sequence is strictly increasing. Anyways, the following is the Zeckendorf decomposition.
\begin{theorem}
    Every nonnegative integer $n$ can be written in exactly one way as the sum of distinct, non-consecutive Fibonacci numbers.
\end{theorem}
We proved this a while ago. The idea behind existence is that the greedy algorithm (take the largest Fiboancci number less or equal to $n$) works. The idea behind uniqueness is that for any set of indices $S$ and $T$ with
\[\sum_{k\in S}F_k=\sum_{\ell\in T}F_\ell\]
must share their largest element because of some bounding; this lets us inductively show $S=T.$ $\blacksquare$

For the sake of analogy, we also have the following similar for powers of $2.$
\begin{proposition}
    Every nonnegative integer $n$ can be written in exactly one way as the sum of distinct powers of $2.$
\end{proposition}
This is, of course, the binary expansion of $2.$ We take this in two parts.
\begin{lemma}
    Every nonnegative integer $n$ can be written in at least one way as the sum of distinct powers of $2.$
\end{lemma}
Like Zeckendorf, we can prove existence by the greedy algorithm. Formally, we do this by strong induction. Our base case is $n=0,$ where no powers of $2$ is one such representation.

Now suppose that all positive integers less than $n>0$ have a binary expansion into distinct powers of $2.$ Because $n\ge1=2^0$ and powers of $2$ grow arbitrarily large, we can find an index $K$ such that
\[2^K\le n<2^{K+1}.\]
(Formally, $K=\floor{\log_2n}.$) Then we use the inductive hypothesis on $n-2^K$ to get a set $S$ so that
\[n-2^K=\sum_{k\in S}2^k\implies n=2^K+\sum_{k\in S}2^k.\]
All terms on the right-hand side are powers of $2,$ so it remains to show that they are all distinct. Well, all elements in $S$ are distinct from each other by hypothesis, so it suffices to show that $K$ is distinct from all elements in $S.$ Well, for any $k\in S,$
\[2^k\le n-2^K<2^{K+1}-2^K=2^K\]
by just stringing our inequalities together. Thus, $K$ is larger than any element in $S,$ implying distinctness. $\blacksquare$

It remains to show uniqueness now. This means that, if $S$ and $T$ are both sets of indices which give decompositions of $n,$ then we must have $S=T.$ In other words, we need
\[\sum_{k\in S}2^k=\sum_{\ell\in T}2^\ell\]
to imply $S=T.$ For this, we have the following lemma which imitates the Zeckendorf case.
\begin{lemma}
    Let $N$ be the largest index in either $S$ or $T.$ Then $N$ must appear in both $S$ and $T.$
\end{lemma}
This comes down to bounding; in some sense, $2^N$ is simply ``too large'' for either $S$ or $T$ to avoid it. Without loss of generality, we'll place $N$ in $S.$ If $N$ is already in $T,$ then we're done. Otherwise, suppose for the sake of contradiction that $N$ isn't in $T.$

Now we bound. Note $T$ is the largest of all indices, so all indices in $T$ are smaller than $N.$ In other words,
\[T\subseteq\{0,2,\ldots,N-1\}.\]
But then
\[\sum_{\ell\in T}2^\ell\le\sum_{\ell=0}^{N-1}2^\ell.\]
Using the geometric sequence sum formula, this is $2^N-1.$ However, $N$ lives in $S,$ so
\[\sum_{k\in S}2^K\ge 2^N>2^N-1\ge\sum_{\ell\in T}2^\ell.\]
This violates the assumption that $S$ and $T$ are both decompositions of the same positive integer! So we have our contradiction. $\blacksquare$

We can now show by induction $\#S$ that we must have $S=T.$ If $\#S=0,$ then we have a sum of $0,$ meaning that $T=\emp$ as well for any elements means positive. Otherwise, suppose the statement is true for $\#S=k,$ and we show $\#S=k+1.$

Well, pick up $S$ and $T$ representing the same positive integer where $\#S=k+1.$ Because $k\ge0,$ there is some element of $S,$ meaning that there is a largest element $N$ among $S$ and $T,$ which by lemma lives in both $S$ and $T.$ This means
\[\sum_{k\in S\setminus\{N\}}2^k=-2^N+\sum_{k\in S}2^k=-2^N+\sum_{\ell\in T}2^\ell=\sum_{\ell\in T\setminus\{N\}}2^\ell.\]
Now, $S\setminus\{N\}$ has size $k$ now, so $S\setminus\{N\}=T\setminus\{N\}.$ Adding $N$ back in, we see $S=T.$ This completes the proof of uniqueness and therefore of the the proposition. $\blacksquare$

The idea here is that powers of $2$ are more or less like playing with the Zeckendorf decomposition ``on easy mode,'' and it's not unreasonable to hope that statements of this flavor which are true for powers of $2$ can be proven in at least a similar way as for Zeckendorf. With that mind, we show the following.
\begin{lemma}
    The first $n$ powers of $2$ can represent exactly the first $2^n$ nonnegative integers.
\end{lemma}
This lemma is essentially saying that there are no ``holes'' in the binary expansion. We can show this quickly after having existence and uniqueness of binary expansions. Namely, any subset $S$ containing only indices for the first $n$ powers of $2$ has $S\subseteq\{0,\ldots,n-1\},$ so
\[\sum_{k\in S}2^k\le\sum_{k=0}^{n-1}2^k=2^n-1\]
because of the geometric sequence sum formula. Thus, binary expansion provides a map
\[S\subseteq\{0,\ldots,n-1\}\longmapsto\sum_{k\in S}2^k\in\left[0,2^n-1\right].\]
Because of uniqueness of representation, this mapping is injective. However, both sides have size $2^n$---the left-hand side is the power set of $\{0,\ldots,n-1\}$ after all---so this mapping must in fact be bijective. The surjective property is the desired statement. $\blacksquare$

The above statement might feel reasonably clear just from looking in binary, but having a formal proof means that we can hope to move it over to the Zeckendorf decomposition as well.
\begin{lemma}
    The first $n$ Fibonacci numbers can represent exactly the first $F_n$ nonnegative integers by Zeckendorf decomposition.
\end{lemma}
We proceed the same way as before but with a bit more care. To start, the largest integer we can represent by Zeckendorf decomposition is
\[F_{n-1}+F_{n-3}+F_{n-5}+\cdots.\]
Indeed, if $S\subseteq\{0,\ldots,n-1\}$ has nonconsecutive terms, then the $k$th largest element of $S$ is no more than $n+1-2k,$ which we can show inductively. (Note $k=1$ is free; then the $k+1$st largest number is at least two less than the $k$th largest.) This gives the above sum.

However, we can actually evaluate this, again inductively. We claim that
\[F_n-1\stackrel?=F_{n-1}+F_{n-3}+F_{n-5}+\cdots=\sum_{k=0}^{\floor{(n-1)/2}}F_{n-1-2k}.\]
When $n=0,$ this reads $F_1-1=1-1=0.$ When $n=1,$ this reads $F_2-1=2-1=1=F_1.$ Then for our inductive step, we suppose $n$ and prove $n+2.$ Indeed,
\[F_{n+2}-1=F_{n+1}+(F_n-1)=F_{n+1}+\sum_{k=0}^{\floor{(n-1)/2}}F_{n-1-2k}=\sum_{k=0}^{\floor{(n+1)/2}}F_{n+1-2k}\]
after shifting indices $k\mapsto k+1.$ This completes the inductive step here.

Thus, we have a map
\[S\subseteq\{0,\ldots,n-1\}\longmapsto\sum_{k\in S}F_k\in[0,F_n-1].\]
Note the domain here is over subsets $S$ with nonconsecutive terms. We know that this mapping is injective, so to show that it is surjective, it suffices to show that the cardinalities of the sets on either side are equal.

Well, the right-hand side is the first $F_n$ nonnegative integers, so we need to show there are $F_n$ subsets of $\{0,\ldots,n-1\}$ with nonconsecutive terms. Well, let $F'_n$ be this number of subsets. We do this by induction. Note that $F'_0=1$ (only the empty set) and $F'_1=2$ ($\{0\}$ and $\emp$).

Then for the inductive step, we fix $n\ge2$ and note that a subset of $\{0,\ldots,n-1\}$ containing nonconsecutive terms either contains $n-1$ or does not. If it contains $n-1,$ then it remains to choose a subset of $\{0,\ldots,n-3\}$ with nonconsecutive terms. If it doesn't contain $n-1,$ then it remains to choose a subset of $\{0,\ldots,n-2\}$ with nonconsecutive terms. Thus,
\[F'_n=F'_{n-1}+F'_{n-2}.\]
Thus, $F'_\bullet$ satisfies the Fibonacci initial conditions and recurrence, so we must have $F'_\bullet=F_\bullet.$ This finishes the proof. $\blacksquare$

That Fibonacci bounding and dealing with nonconsecutive terms is exactly what makes the Zeckendorf decomposition ``hard mode'' here, but the overall outline is still intact. Anyways, our lemma quickly gives us the following: the unique decomposition property implies power of two.
\begin{proposition}
    Suppose $\{a_k\}_{k=0}^\infty$ is an increasing sequence of positive integers for which every nonnegative integer $n$ has a unique subset $S\subseteq\NN$ such that
    \[\sum_{k\in S}a_k=n.\]
    Then $a_k=2^k$ for each $k\in\NN.$
\end{proposition}
We do this by induction. Note that $n=0$ has $S=\emp,$ but $n=1$ needs to have some nonempty subset $S_1.$ This means that
\[a_0\le\sum_{k\in S_1}a_k=1.\]
Because $a_0$ is a positive integer, we know $a_0\ge1,$ so we conclude $a_0=1=2^0.$ This is our base case.

Now assume that each $k<n$ has $a_k=2^k.$ We show that $a_n=2^n$ for our inductive step. Well by lemma, the first $n$ terms $a_0,\ldots,a_{n-1}$ (which are powers of $2$) represent exactly the first $2^n$ nonnegative integers. We will get what we want by asking how to represent $2^n.$ Let $S$ represent $2^n.$

Note that $S$ cannot be a subset of $\{0,\ldots,n-1\},$ for subsets of these indices represent the nonnegative integers less than $2^n.$ So, letting $m\in S$ have $m\ge n,$ we see that
\[a_n\le a_m\le\sum_{k\in S}a_k=2^n.\]
On the other hand, every nonnegative integer less than $2^n$ already has a representation by indices $\{0,\ldots,n-1\},$ so we must have $a_n\ge2^n$ due to uniqueness---otherwise ``$a_n$'' is a valid decomposition. It follows that $a_n=2^n.$ This competes the proof. $\blacksquare$

Of course, we wouldn't have given this proof if it did not also carry over to the Zeckendorf case.
\begin{theorem}
    Suppose $\{a_k\}_{k=0}^\infty$ is an increasing sequence of positive integers for which every nonnegative integer $n$ has a unique subset $S\subseteq\NN$ with nonconsecutive terms such that
    \[\sum_{k\in S}a_k=n.\]
    Then $a_k=F_k$ for each $k\in\NN.$
\end{theorem}
Again, we induct. As before, $n=0$ has $S=\emp,$ but $n=1$ needs to be represented by a nonempty subset $S_1.$ Then
\[a_0\le\sum_{k\in S}a_k=1.\]
Because $a_0$ needs to be positive, we conclude $a_0=1=F_0.$

Now we assume that each $k<n$ has $a_k=F_k$ and will show $a_n=F_n.$ Our lemma this time says that the first $n$ terms $a_0,\ldots,a_{n-1}$ represent exactly the nonnegative integers less than $F_n,$ so we ask how to represent $F_n.$ Let $S$ be the indices for $F_n.$

As before, we can lower-bound $a_n$ by $F_n$ because all positive integers less than $F_n$ all have Zeckendorf decompositions: having $a_n<F_n$ would give $a_n$ two decompositions as the one we had before as well as ``$a_n$.'' To upper-bound, we note $S$ representing $F_n$ needs an index at least $n$ because
\[S\subseteq\{0,\ldots,n-1\}\implies\sum_{k\in S}F_k<F_n.\]
Thus,
\[a_n\le a_m\le\sum_{k\in S}a_k=F_n.\]
Combining inequalities, we conclude $a_n=F_n,$ which completes the induction and so the theorem. $\blacksquare$

This is sufficiently cute, so we call it here. As some end commentary, my entries here are about to becomes much shorter because PROMYS is starting up, and this is no longer my full-time project.

\subsubsection{July 2nd}
Today I learned that integrals with $\pi(x)$ can actually have closed forms, from \href{https://math.stackexchange.com/a/2270227/869257}{here}. Namely, I'm used to being forced to using $\pi(x)\sim\frac x{\log x}$ or Abel summation to get reasonable bounds. Here is the integral of interest.
\[I:=\int_0^\infty\frac{\pi(x)}{x\left(x^s-1\right)}\,dx=\int_2^\infty\frac{\pi(x)}{x\left(x^s-1\right)}\,dx.\]
Very quickly, showing that this converges is not difficult.
\begin{lemma}
    The integral $I$ converges for $\op{Re}(s)>1.$
\end{lemma}
We bound stupidly. Note $\pi(x)\le x$ gives
\[|I|\le\left|\int_2^\infty\frac x{x\left(x^s-1\right)}\,dx\right|=\left|\int_2^\infty\frac1{x^s-1}\,dx\right|.\]
Using the triangle inequality, $\left|x^s-1\right|\le\left|x^s\right|-1,$ so we can bound
\[|I|\le\int_2^\infty\frac1{x^{\op{Re}(s)}-1}\,dx,\]
where $\left|x^s\right|=x^{\op{Re}(s)}$ for real $x>0.$ To bound this cleanly, we note that $\op{Re}s>1$ gives $x^{\op{Re}s}\ge2^{\op{Re}s}>2$ over this interval, so
\[\frac1{x^{\op{Re}s}-1}\le\frac2{x^{\op{Re}s}}\]
after some rearranging. Thus,
\[|I|\le\int_2^\infty\frac2{x^{\op{Re}s}}\,dx=\frac{2x^{1-\op{Re}s}}{1-\op{Re}s}\bigg|_2^\infty=\frac{2\cdot2^{1-\op{Re}s}}{\op{Re}s-1}.\]
This is indeed finite, so we are done. $\blacksquare$

I think it's helpful to watch the evaluation expecting certain things to come together, s   o we will go ahead and give away the answer now.
\begin{proposition}
    We can evaluate, for $\op{Re}s>1,$
    \[I=\frac1s\log\zeta(s).\]
\end{proposition}
Before going into the evaluation, we say at the outset that what makes this integral tick is that
\[\frac d{dx}\log\left(1-\frac1{x^s}\right)=\frac1{1-\frac1{x^s}}\cdot\frac{-s}{x^{s+1}}=\frac{-s}{x\left(x^s-1\right)}.\]
This will come into play shortly.

Because we are actually trying to evaluate $I,$ the cleanest thing to do is split up by primes so that $\pi(x)$ is well-behaved. Because we already know that the integral converges, it suffices to evaluate the integral on a subsequence of upper bounds. We write
\[I=\int_2^\infty\frac{\pi(x)}{x\left(x^s-1\right)}\,dx=\sum_{k=1}^\infty\int_{p_k}^{p_{k+1}}\frac k{x\left(x^s-1\right)}\,dx,\]
where $p_k$ is the $k$th prime. (To be clear, the subsequence of upper bounds we are using are the $p_\bullet.$) Now $k$ is a constant, so we can use the derivative we found earlier to see that the inner integral evaluates as
\[I=\sum_{k=1}^\infty k\cdot-\frac1s\cdot\log\left(1-\frac1{x^s}\right)\bigg|_{p_k}^{p_{k+1}}.\]
Simplyifying slightly, we have
\[sI=-\sum_{k=1}^\infty k\left[\log\left(1-\frac1{p_k^s}\right)-\log\left(1-\frac1{p_{k+1}^s}\right)\right].\]
At this point, we can see that we are going to get $\zeta$ from the Euler product. We would like this sum to telescope to hopefully rid of the $k$ on the outside, and even though it is possible to manipulate the infinite sums directly, we will go ahead fix a $N$ as the upper bound and manipulate like that. This looks like
\[(sI)_N=-\sum_{k=1}^Nk\log\left(1-\frac1{p_k^s}\right)+\sum_{k=1}^Nk\left(1-\frac1{p_{k+1}^s}\right).\]
In the second sum, we can shift $k\mapsto k-1$ to telescope. Note that we are allowed to keep the lower index as $k=1$ because we end up multiplying the term by $(k-1),$ which vanishes. This is
\[(sI)_N=-\sum_{k=1}^Nk\log\left(1-\frac1{p_k^s}\right)+\sum_{k=1}^{N+1}(k-1)\log\left(1-\frac1{p_k^s}\right).\]
Now we see the telescoping, giving
\[(sI)_N=-\sum_{k=1}^N\log\left(1-\frac1{p_k^s}\right)+N\log\left(1-\frac1{p_{N+1}^s}\right).\]
We hope that the remainder term will vanish as $N\to\infty,$ but let's focus on the main term right now. Moving the logarithm out into a product, we see it is
\[-\sum_{k=1}^N\log\left(1-\frac1{p_k^s}\right)=\log\prod_{k=1}^N\frac1{1-p_k^{-s}}.\]
We see that $N\to\infty$ makes this the Euler product for $\zeta(s).$ Thus, the main term is indeed giving the desired evaluation.

So indeed, it remains to show that $N\log\left(1-p_{N+1}^{-s}\right)$ actually vanishes as $N\to\infty.$
\begin{lemma}
    For $\op{Re}s>1,$ we have that $N\log\left(1-p_{N+1}^{-s}\right)\to0$ as $N\to\infty.$
\end{lemma}
Being careful with the $s\in\CC,$ we expand the logarithm directly as
\[\log\left(1-\frac1{p_{N+1}^s}\right)=-\sum_{k=1}^\infty\frac1{kp_{N+1}^{ks}}.\]
This means that
\[\left|\log\left(1-\frac1{p_{N+1}^s}\right)\right|\le\sum_{k=1}^\infty\frac1{kp_{N+1}^{k\op{Re}s}}=-\log\left(1-\frac1{p_{N+1}^{\op{Re}s}}\right).\]
So we let $\sigma:=\op{Re}s>1$ and forget about $s.$ Now, because $N<p_{N+1},$ we see
\[\left|N\log\left(1-\frac1{p_{N+1}^\sigma}\right)\right|<\left|p_{N+1}\log\left(1-\frac1{p_{N+1}^\sigma}\right)\right|.\]
So let $x:=p_{N+1}^{-1}$ and investigate $x^{-1}\log\left(1-x^\sigma\right)$ as $x\to0.$ From here, L'H\^opital's Rule can finish by writing
\[\lim_{x\to 0}\frac{\log\left(1-x^\sigma\right)}x.\]
As $x\to0,$ both the numerator and denominator approach $0,$ so we will apply L'H\^opital's. The derivative of the numerator is $\left(1-x^\sigma\right)^{-1}\cdot\sigma x^{\sigma-1},$ so this limit is
\[\lim_{x\to0}\left(1-x^\sigma\right)^{-1}\cdot\sigma x^{\sigma-1}=(1-0)^{-1}\cdot\sigma0^{\sigma-1}=0.\]
(Here is where $\sigma>1$ matters!) This finishes the proof. $\blacksquare$

Combining everything we've done, we see that
\[(sI)_N=-\sum_{k=1}^N\log\left(1-\frac1{p_k^s}\right)+N\log\left(1-\frac1{p_{N+1}^s}\right).\]
As $N\to\infty,$ the sum becomes $\log\zeta(s),$ and the remainder term vanishes, so we do get
\[sI=\log\zeta(s).\]
This rearranges into the desired result. $\blacksquare$

What I like about this result is that it almost feels like it could be rederivable directly from trying to use
\[\int\frac{-s}{x\left(x^s-1\right)}\,dx=\log\left(1-\frac1{x^s}\right)+C.\]
Namely, summing this over all primes will give us $\log\zeta(s)$ on the right-hand side if we deal with $C$ properly, and it feels like the telescoping is done entirely to make this $C$ dealt with.

Granted, the telescoping part feels a bit unmotivated, but I can't shake the feeling that it's really just integration by parts/summation by parts with $\pi(x).$ I haven't been able to formalize this intuition, but I would not be surprised if a simplification of this form were possible.

\subsubsection{July 3rd}
Today I learned that there are no three-dimensional division algebras over the reals. \todo{}
